\documentclass[phy1200]{handout}

\title{Blue sky}

\begin{document}

\maketitle

An explanation of why the sky is blue has two pieces: you have to understand sunlight and you have to understand what light does when it hits small particles.

You and the group across from you are both using \textbf{spectroscopes}. These let you image a small beam of light and see what wavelengths are inside of it. Light coming into a spectroscope hits something (usually a prism or diffraction grating) and the different wavelengths leave the thing at different angles. Consider how a prism might break light up into various wavelengths using refraction.

Once you have a spectroscope, you can look at the \textbf{spectrum} of a light source: its wavelength fingerprint. Open up `spectroscopy', a program on your computer. It won't recognize your spectroscope right away, you need to open up the end. This receives a specific beam of light to send into the spectroscope. Once it recognizes the spectroscope, you can hit `analyze light' at the top and start recording spectra down in the bottom left.

Try pointing your spectroscope at the optics source (an incandescent bulb) and see what spectrum is measured. Then try your phone camera. If you want to try the sun, let the instructor know and you can use a laptop with the software. We also have a color source with red, blue, green, and white LEDs. You can look at those in the spectroscope too. Try to wrap your head around how we our eyes interpret a spectrum as color.

\section{Blackbody spectra}

Go read a little bit about blackbody spectra.\footnote{Open Stax university physics, volume 3 chapter 6.1.} Blackbody spectra are emitted by anything above absolute zero. This includes incandescent bulbs and the sun. Even you are emitting blackbody radiation, but in the infrared! The peak wavelength is related to temperature by a formula called Wien's (pronounced ``Veen'') displacement law. Find that formula and explain what it means. 

What temperatures would peak inside the visible spectrum? Why don't we see most humans glowing? Measure the spectrum of the bulb in the optics source and figure out a rough estimate of its temperature. Measure sunlight and figure out a rough estimate of the temperature at the surface of the sun.

Your phone camera isn't a great match for a blackbody source (it is actually an LED, which emits more focused colors without heating up), but figure out the approximate temperature of the light using Wien's law. It is common to report temperature (in K) like this for light bulbs, look around next time you go to the store.\footnote{Does hotter or colder temperature correspond to a `warm' or `cool' color?}

\section{Blue sky and scattering}
Go get some water from the sink (there is one in DSC 053) and fill the clear tank. This is currently transparent over short distances. So is air. However, over long distance, light scatters off of air molecules and dust. We can speed this process up to get a whole sunset in a tank. Mix in small amounts of milk to the tank and shine your phone light into the tank. Keep adding small amounts of milk until you start to get orange-ish light at the far end of the tank.

Measure the spectrum of the scattered light by holding your spectroscope up to the tank and scanning it down the tank. There is a wavelength dependence to what light is preferentially scattered (Rayleigh scattering). What do you notice about which colors get scattered more strongly? How is the beam changing as it goes through the tank? Why is it orange toward the end? 

Look at your phone light through the far side of the tank. This is directly analogous to the view of a sunset. Explain why the light is so red now.

\section{Report}
In your report be sure to explain:
\begin{enumerate}
    \item What a spectroscope does, what a spectrum is.
    \item What a blackbody spectrum is, how it changes with temperature.
    \item How to find a temperature from a blackbody spectrum, ballpark estimates for the temperature of the bulb/sun.
    \item Explain the wavelength dependence of scattering (for small particles).
    \item Explain how this leads to the sky being blue and the sunset being red.
\end{enumerate}







\end{document}